% Generated draft: requires mathematical and source review.
\documentclass[11pt]{article}
\usepackage[margin=25mm]{geometry}
\usepackage{fontspec}
\setmainfont{Latin Modern Roman}
\usepackage{amsmath,amssymb,mathtools,unicode-math}
\setmathfont{Latin Modern Math}
\usepackage{enumitem,needspace,xurl,hyperref}
\setlength{\parindent}{0pt}
\setlength{\parskip}{5pt}
\setlength{\emergencystretch}{3em}
\newcommand{\N}{\mathbb{N}}
\newcommand{\Z}{\mathbb{Z}}
\newcommand{\Q}{\mathbb{Q}}
\newcommand{\R}{\mathbb{R}}
\newcommand{\C}{\mathbb{C}}
\newcommand{\F}{\mathbb{F}}
\newcommand{\A}{\mathbb{A}}
\newcommand{\PP}{\mathbb{P}}
\newcommand{\E}{\mathbb{E}}
\newcommand{\eps}{\varepsilon}
\newcommand{\dd}{\,\mathrm{d}}
\begin{document}
\section*{UnsolvedMath ID: 701268}
\textbf{Source problem code:} catalogue entry
\subsection{Definitions and mathematical statement}
An involution in a group is an element of order exactly \(2\). A group \(H\) is called \emph{strictly real} if
\[
\forall h\in H\setminus\{1\}\quad \exists t\in H\quad
\bigl(t\ne 1,\ t^2=1,\ tht^{-1}=h^{-1}\bigr).
\]
The involution may depend on \(h\).

Let \(k\) be a field of characteristic \(2\), and let \(G\) be a group of Lie type over \(k\). Write \(\mathcal U(G)\) for its maximal unipotent subgroups, understood in the intended Lie-type setting. Determine which such groups \(G\) satisfy
\[
\forall U\in\mathcal U(G)\quad
\forall u\in U\setminus\{1\}\quad
\exists t\in U\quad
\bigl(t\ne 1,\ t^2=1,\ tut^{-1}=u^{-1}\bigr).
\]
In particular, the conjugating involution must belong to \(U\) itself.

The source does not specify whether the field is finite, which forms of groups of Lie type are included, or the precise setting in which maximal unipotent subgroups are taken. These scope qualifications remain unresolved; the displayed question retains them rather than restricting the problem to a particular class.
\subsection{Short English statement}
Classify the groups of Lie type over fields of characteristic \(2\) whose maximal unipotent subgroups are strictly real: every nonidentity element of each such subgroup must be inverted by conjugation by an involution within that subgroup. The intended scope of groups and fields requires clarification.
\subsection{Sources}
\begin{itemize}
\item Local UnsolvedMath registry: permanent ID \(701268\), exact source problem code \texttt{catalogue entry}, title \emph{Kourovka Notebook Problem 16.76}.
\item Supplied attribution: Ya. N. Nuzhin, proposed in \(2006\); \emph{The Kourovka Notebook (Unsolved Problems in Group Theory)}, 16th issue (2006), Problem 16.76.
\item The supplied background also identifies the 21st issue edition (2026), \texttt{arXiv:1401.0300v48}. The collection description names E. I. Khukhro and V. D. Mazurov as editors. These bibliographic details are supplied source data and have not been independently verified.
\item The supplied registry labels the problem \emph{open}. This is a supplied status label, not an independently verified assessment of the literature.
\end{itemize}
\end{document}
